\documentclass[11pt]{article}

\usepackage[margin=1in]{geometry}
\usepackage[T1]{fontenc}
\usepackage[utf8]{inputenc}
\usepackage{lmodern}
\usepackage{microtype}
\usepackage{amsmath, amssymb, amsthm, mathtools}
\usepackage{booktabs}
\usepackage{enumitem}
\usepackage{xcolor}
\usepackage[colorlinks=true, linkcolor=blue!55!black, citecolor=blue!55!black, urlcolor=blue!55!black]{hyperref}

% -----------------------------------------------------------------------------
% Notation macros
% -----------------------------------------------------------------------------

\newcommand{\R}{\mathbb{R}}
\newcommand{\E}{\mathbb{E}}
\newcommand{\ip}[2]{\left\langle #1, #2 \right\rangle}
\newcommand{\norm}[1]{\left\lVert #1 \right\rVert}
\newcommand{\softmax}{\operatorname{softmax}}
\newcommand{\diag}{\operatorname{diag}}

\newcommand{\dm}{d_{\mathrm{model}}}
\newcommand{\dhead}{d_{\mathrm{head}}}
\newcommand{\dsae}{d_{\mathrm{SAE}}}

\newcommand{\pretag}{\mathrm{pre}}
\newcommand{\midtag}{\mathrm{mid}}
\newcommand{\attn}{\mathrm{attn}}
\newcommand{\OV}{\mathrm{OV}}
\newcommand{\QK}{\mathrm{QK}}
\newcommand{\LN}{\mathrm{LN}}

\newcommand{\hook}[1]{\texttt{#1}}
\newcommand{\residpre}{\mathbf{r}^{\pretag}}
\newcommand{\residmid}{\mathbf{r}^{\midtag}}
\newcommand{\attnout}{\mathbf{attn\_out}}

\newcommand{\todo}[1]{\textcolor{red!65!black}{\textbf{TODO:} #1}}

% -----------------------------------------------------------------------------
% Theorem-like environments
% -----------------------------------------------------------------------------

\theoremstyle{definition}
\newtheorem{definition}{Definition}
\newtheorem{protocol}{Protocol}

\theoremstyle{plain}
\newtheorem{proposition}{Proposition}
\newtheorem{claim}{Claim}

\theoremstyle{remark}
\newtheorem{remark}{Remark}

% -----------------------------------------------------------------------------

\title{Feature-Resolved Attention: A Clean Derivation}
\author{Dmitry Manning-Coe}
\date{\today}

\begin{document}

\maketitle

\begin{abstract}
Feature-resolved attention (FRA) decomposes an attention circuit in the basis
of sparse autoencoder (SAE) features.  The goal is to assign a scalar target,
such as a downstream feature preactivation or steering direction, to specific
query, key, and value-side feature contributions.  This note gives a clean
derivation of the two pieces we use most often: frozen-attention OV attribution
and first-order QK attribution via the softmax Jacobian.  It also records the
distinction between feature ranking and intervention path, since conflating
these two choices has been a recurring source of confusion.
\end{abstract}

\tableofcontents

% =============================================================================
\section{Decomposition}
% =============================================================================
Our goal is to decompose the attention block into linearly independent contributions from each feature. To do so, we must first obtain the decomposition of the residual stream vector $\mathbf{x}$ into a linear combination $\mathbf{u}$ of Sparse Autoencoder features:
\begin{equation}
    \mathbf{u}\equiv\sigma\left(W^{\textrm{enc}}\mathbf{x}+\mathbf{b}^{\textrm{enc}}\right).
\end{equation}
We will typically work in component form and so we note that:
\begin{equation}
    u^\lambda_{q}=\sigma\left(\sum_{a}
       W^{\textrm{enc},\lambda}_a x_{q a}+b^{\textrm{enc},\lambda}
    \right).
\end{equation}
Note that in the SAE literature it is conventional to associate each latent feature $\lambda$ with a decoder direction $\mathbf{f}^\lambda$ in residual-stream space. Under the column-vector convention used here, this is the image of the latent basis vector $\mathbf{e}^\lambda$ under the decoder:
\begin{equation}
    \mathbf{f}^\lambda\equiv W^{\textrm{dec}}\mathbf{e}^\lambda
    = W^{\textrm{dec}}_{:,\lambda}.
\end{equation}
Equivalently, in row-vector SAE implementations where $\mathbf{x}'=\mathbf{u}W^{\textrm{dec}}+\mathbf{b}^{\textrm{dec}}$, the same feature direction is the corresponding row $W^{\textrm{dec}}_{\lambda:}$.

The latent vector $\mathbf{u}$ is then projected back by the decoder matrix to obtain the reconstructed activation:
\begin{equation}
    \mathbf{x}'\equiv W^{dec}\mathbf{u}+\mathbf{b}^{dec}.
\end{equation}
This reconstruction will not, in general, perfectly recover the original activation. The difference defines the reconstruction error $\boldsymbol{\varepsilon}$\footnote{Note that, since we do know the exact activations in all the cases we consider, we can measure the error term and see how it interacts with the features!}:
\begin{equation}
    \mathbf{x}=\mathbf{x}'+\boldsymbol{\varepsilon}.
\end{equation}
The component-wise attention update from the resid-pre hookpoint is given by:
\begin{equation}
    x^{\textrm{mid}}_{q a}=x^{\textrm{pre}}_{q a}+Attn[\mathbf{x}^{\textrm{ln1}}]_{q a}.
\end{equation}
To lighten notation, we use the convention that $\mathbf{x}$ with no superscript is the activation vector at the $ln1$ hookpoint. Following \cite{elhage_mathematical_framework}, we decompose the attention block into an attention pattern $A$ and an $OV$ component:
\begin{equation}
    Attn(\mathbf{x})_{q a}=\sum_{k}A(\mathbf{x})_{qk}OV(\mathbf{x})_{k a}.
\end{equation}
The attention pattern $A$ is given by the contraction of the $Q$ and $K$ matrices for each head $h$:
\begin{equation}
    A(\mathbf{x})_{qk}=\frac{1}{\sqrt{d_h}}\sum_{h}\sum_{a,b,i}
    \left(
        Q^h_{i a}x_{q a}+b^{Q,h}_{i}
    \right)
    \left(
        K^h_{i b}x_{k b}+b^{K,h}_{i}
    \right),
\end{equation}
expanding out the terms:
\begin{equation}
    A(\mathbf{x})_{qk}=\frac{1}{\sqrt{d_h}}\sum_{h}\sum_{a,b,i}
    Q^h_{i a}x_{q a}K^h_{i b}x_{k b}
    +b^{Q,h}_{i}K^h_{i b}x_{k b}
    +Q^h_{i a}x_{q a}b^{K,h}_{i}
    +b^{Q,h}_{i}b^{K,h}_{i}
\end{equation}
This is a bit of a mess, so we can probably do better by dropping component form and writing:
\begin{equation}
    A(\mathbf{x})=\mathbf{x}^TQ^TK\mathbf{x}+\underbrace{(\mathbf{b}^Q)^TK}_{\equiv (\mathbf{B}^Q)^T}\mathbf{x}+\mathbf{x}^T\underbrace{Q^T\mathbf{b}^K}_{\equiv \mathbf{B}^K}+\underbrace{(\mathbf{b}^Q)^T\mathbf{b}^K}_{\equiv B^{QK}}.
\end{equation}
The $OV$ component is given as:
\begin{equation}
    OV(\mathbf{x})_{q a}=\sum_h\sum_{a,i}
    \left(
        W^{O,h}_{a i}\left(W^{V,h}_{i b}x_{q b}+b^O_{i}\right)+b^{V,h}_{a}
    \right),
\end{equation}
or, in vector form:
\begin{equation}
   OV(\mathbf{x})=W^O(W^V\mathbf{x}+\mathbf{b}^V)+\mathbf{b}^O=W^OW^V\mathbf{x}+\mathbf{B}^{OV},
\end{equation}
where we define:
\begin{equation}
    \mathbf{B}^{OV}\equiv W^O\mathbf{b}^V+\mathbf{b}^O.
\end{equation}
To feature-resolve, we first need to insert the reconstructed activations $\mathbf{x}=\mathbf{x}'+\boldsymbol{\varepsilon}$. For completeness, we keep track of the error term explicitly. The attention pattern becomes:
\begin{equation}
A(\mathbf{x}'+\boldsymbol{\varepsilon})=A(\mathbf{x}')+A(\boldsymbol{\varepsilon})+\boldsymbol{\varepsilon}^TQ^TK\mathbf{x}'+(\mathbf{x}')^TQ^TK\boldsymbol{\varepsilon}-B^{QK},
\end{equation}
where the last term just comes from overcounting the bias. The $OV$ term is:
\begin{equation}
    OV(\mathbf{x}'+\boldsymbol{\varepsilon})=OV(\mathbf{x}')+OV(\boldsymbol{\varepsilon})-\mathbf{B}^{OV}.
\end{equation}
To feature-resolve, we simply propagate through the feature-wise decomposition of the activations. That is, we replace:
\begin{equation}
x'_{q a}\to X^\lambda_{q a}\equiv W^{dec,\lambda}_{a}u^\lambda_{q}+b^{dec}_{a},
\end{equation}
through the attention block. We note that the reconstruction error $\boldsymbol{\varepsilon}$ is defined only after reconstruction, and hence cannot be decomposed into underlying features. Putting the feature-resolved activation $X$ into the attention pattern gives:
\begin{equation}
    A(\mathbf{x}')=XQ^TKX+(\mathbf{B}^Q)^TX+X^TQ^T\mathbf{B}^K+B^{QK}
\end{equation}
In component form:
\begin{equation}
    A(\mathbf{x}')_{q k}=\sum_{h=1}^{n}\sum_{i=1}^{d_h}\sum_{a,b=1}^{d_M}X^{\mu}_{q a}Q_{a i}K_{i b}X^\nu_{b k}+B^Q_{a} X^\nu_{a q}+X^\mu_{q a}(B^K)_{a}+B^{QK},
\end{equation}
and similarly for the attention pattern on $\boldsymbol{\varepsilon}$. Doing the same for the $OV$:
\begin{equation}
    OV(\mathbf{x}')=W^{OV}X+\mathbf{B}^{OV},
\end{equation}
becomes in component form:
\begin{equation}
    OV(\mathbf{x}')^{\lambda}_{q a}=\sum_{}W^{OV}_{a b}X^\lambda_{q b}+B^{OV}_{a}.
\end{equation}
Putting all of this together, we can decompose the full attention block as:
\begin{align}
    Attn(x)_{q a}&=\sum_{k}A(x'+\varepsilon)_{q k}OV(x'+\varepsilon)_{k a}\\
    &=\sum_{k}\left[        A(\mathbf{x}')+A(\boldsymbol{\varepsilon})+\boldsymbol{\varepsilon}^TQ^TK\mathbf{x}'+(\mathbf{x}')^TQ^TK\boldsymbol{\varepsilon}-B^{QK}
    \right]_{qk}
    \left[
    OV(\mathbf{x}')+OV(\boldsymbol{\varepsilon})-\mathbf{B}^{OV}
    \right]_{k a}
\end{align}
We note that models that use Rotary Positional Embeddings (RoPE) can be accommodated in this framework by introducing modified position-dependent $QK$ matrices:
\begin{equation}
    \tilde Q^h_{q,ia}=R^h_{q,ij}Q^h_{ja},
  \qquad
  \tilde K^h_{k,ib}=R^h_{k,ij}K^h_{jb},
\end{equation}


Schematically: 
Standard:
\begin{equation}
    A(\mathbf{x})_{qk}\sim \sum_{a,b,i}
    \left(
        Q_{i a}x_{q a}+b^{Q}_{i}
    \right)
    \left(
        K_{i b}x_{k b}+b^{K}_{i}
    \right),
\end{equation}
FRA:
\begin{equation}
    A^{\mu\nu}(\mathbf{x})_{qk}\sim \sum_{a,b,i}
    \left(
        Q_{i a}X^\mu_{q a}+b^{Q}_{i}
    \right)
    \left(
        K_{i b}X^\nu_{k b}+b^{K}_{i}
    \right),
\end{equation}

OV:
\begin{equation}
    OV_{q a}=\sum_{b}OV_{a b}x_{q b}
\end{equation}
FRA-OV:

\begin{equation}
    OV^\lambda_{q a}=\sum_{b}OV^\lambda_{a b}X^\lambda_{q b}
\end{equation}







% =============================================================================
\appendix
\section{Appendix Placeholders}
% =============================================================================

% \subsection{LayerNorm linearization}
% \todo{Insert the resid-pre to ln1 two-stage derivation if this writeup needs
% the sleeper-specific two-stage path.}

% \subsection{Channel-wise interaction decomposition}
% \todo{Insert the Q/K/V coalition or Harsanyi decomposition if the paper needs
% singles, pairs, and triples rather than the frozen-OV/QK split.}

% \subsection{Experimental protocol}
% \todo{Insert exact prompts, feature counts, alpha sweeps, and metrics for each
% model family.}

\end{document}
